\documentclass{article}
\usepackage{booktabs} 
\setlength{\floatsep}{5pt}      
\setlength{\textfloatsep}{10pt}
\setlength{\intextsep}{5pt}   
\usepackage{forest}
\usepackage{float}

% Definizione dello stile per gli alberi di ricerca
\forestset{
  explored/.style={draw=olive!80!black, thick, fill=yellow!20},
  unexplored/.style={draw=gray, thick, fill=white}
}

\begin{document}
\section*{Esplorazione Spazio degli Stati}

% ---------------------------------------------------------
% Tabella DFS
% ---------------------------------------------------------
\begin{table}[H]
    \centering
    \caption{ricerca in profondità (DFS)}
    \begin{tabular}{clp{4cm}p{5cm}}
        \toprule
        \textbf{Passo} & \textbf{Nodo Espanso} & \textbf{Frontiera} & \textbf{Esplorati} \\
        \midrule
        1 & S  & $\emptyset$    & $\emptyset$ \\
        2 & A  & [B, D]         & \{S\} \\
        3 & B  & [D]            & \{A, S\} \\
        4 & C  & [D, G1]        & \{B, A, S\} \\
        5 & F  & [D, G1, J]     & \{B, C, A, S\} \\
        6 & D  & [G1, J, G2]    & \{F, B, C, A, S\} \\
        7 & E  & [G1, J, G2]    & \{F, B, D, C, A, S\} \\
        8 & G2 & [G1, J]        & \{E, F, B, D, C, A, S\} \\
        \bottomrule
    \end{tabular}
    
    \vspace{0.2cm}
    \textbf{Cammino:} S $\rightarrow$ A $\rightarrow$ B $\rightarrow$ C $\rightarrow$ F $\rightarrow$ D $\rightarrow$ E $\rightarrow$ G2 \quad \textbf{Costo:} 22

\begin{center}
\begin{forest}
  for tree={circle, minimum size=1.5em, inner sep=2pt, s sep=4mm, l sep=5mm, unexplored}
  [S, explored
    [A, explored
      [B, explored
        [C, explored
          [F, explored
            [D, explored
              [E, explored
                [G2, explored]
              ]
            ]
            [G2]
          ]
          [J]
        ]
        [G1]
      ]
    ]
    [B]
    [D]
  ]
\end{forest}
\end{center}
\end{table}
\begin{itemize}
    \item sono state necessarie 8 iterazioni (nodi espansi) per trovare l'obiettivo
    \item il cammino trovato non è ottimo (22 $>$ 14).
    \item la DFS non garantisce l'ottimalità, poiché esplora i rami in profondità senza considerare i costi dei passi
\end{itemize}


% ---------------------------------------------------------
% Tabella BFS
% ---------------------------------------------------------
\begin{table}[H]
    \centering
    \caption{ricerca in ampiezza (BFS)}
    \begin{tabular}{clp{4cm}p{5cm}}
        \toprule
        \textbf{Passo} & \textbf{Nodo Espanso} & \textbf{Frontiera} & \textbf{Esplorati} \\
        \midrule
        1 & S  & $\emptyset$    & $\emptyset$ \\
        2 & A  & [B, D]         & \{S\} \\
        3 & B  & [D]            & \{A, S\} \\
        4 & D  & [C, G1]        & \{B, A, S\} \\
        5 & C  & [G1, E]        & \{B, D, A, S\} \\
        6 & E  & [G1, F, J]     & \{B, C, D, A, S\} \\
        7 & G1 & [F, J, G2]     & \{E, B, C, D, A, S\} \\
        \bottomrule
    \end{tabular}
    
    \vspace{0.2cm}
    \textbf{Cammino:} S $\rightarrow$ B $\rightarrow$ G1 \quad \textbf{Costo:} 16

\begin{itemize}
    \item sono state necessarie 7 iterazioni per trovare l'obiettivo    
    \item il cammino trovato non è ottimo $16 > 14$
    \item BFS garantisce l'ottimalità solo se i costi degli archi sono costanti (costo di cammino è funzione monotona della profondità)
\end{itemize}

\begin{center}
\begin{forest}
  for tree={circle, minimum size=1.5em, inner sep=2pt, s sep=4mm, l sep=5mm, unexplored}
  [S, explored
    [A, explored]
    [B, explored
      [C, explored]
      [G1, explored]
    ]
    [D, explored
      [E, explored]
    ]
  ]
\end{forest}
\end{center}
\end{table}

% ---------------------------------------------------------
% Tabella UCS
% ---------------------------------------------------------
\begin{table}[H]
    \centering
    \caption{ricerca a costi uniformi (UCS)}
    \begin{tabular}{clp{4cm}p{5cm}}
        \toprule
        \textbf{Passo} & \textbf{Nodo Espanso} & \textbf{Frontiera} & \textbf{Esplorati} \\
        \midrule
        1 & S  & $\emptyset$    & $\emptyset$ \\
        2 & A  & [B, D]         & \{S\} \\
        3 & D  & [B]            & \{A, S\} \\
        4 & B  & [C, E]         & \{D, A, S\} \\
        5 & C  & [E, G1]        & \{B, D, A, S\} \\
        6 & E  & [G1, F, J]     & \{B, C, D, A, S\} \\
        7 & F  & [G1, J, G2]    & \{E, B, C, D, A, S\} \\
        8 & J  & [G1, G2]       & \{F, E, B, C, D, A, S\} \\
        9 & G2 & [G1]           & \{J, F, E, B, C, D, A, S\} \\
        \bottomrule
    \end{tabular}
    
    \vspace{0.2cm}
    \textbf{Cammino:} S $\rightarrow$ D $\rightarrow$ C $\rightarrow$ F $\rightarrow$ G2 \quad \textbf{Costo:} 14

\begin{center}
\begin{forest}
  for tree={circle, minimum size=1.5em, inner sep=2pt, s sep=4mm, l sep=5mm, unexplored}
  [S, explored
    [A, explored
      [B, explored
        [G1]
      ]
    ]
    [D, explored
      [C, explored
        [F, explored
          [G2, explored]
        ]
        [J, explored]
      ]
      [E, explored]
    ]
  ]
\end{forest}
\end{center}
\end{table}

\begin{itemize}
    \item sono state necessarie 9 iterazioni per trovare l'obiettivo
    \item il cammino trovato è ottimo (costo 14)
    \item UCS garantisce sempre l'ottimalità se i costi degli archi sono non negativi (e l'albero di ricerca è finito)
\end{itemize}

% ---------------------------------------------------------
% Tabella IDS
% ---------------------------------------------------------
\begin{table}[H]
    \centering
    \caption{ricerca ad approfondimento iterativo (IDS)}
    \begin{tabular}{cclp{3cm}p{4cm}}
        \toprule
        \textbf{Limite} & \textbf{Passo} & \textbf{Nodo Espanso} & \textbf{Frontiera} & \textbf{Esplorati} \\
        \midrule
        \textbf{0} & 1 & S  & $\emptyset$ & $\emptyset$ \\
        \midrule
        \textbf{1} & 1 & S  & $\emptyset$ & $\emptyset$ \\
                   & 2 & A  & [B, D]      & \{S\} \\
                   & 3 & B  & [D]         & \{A, S\} \\
                   & 4 & D  & $\emptyset$ & \{B, A, S\} \\
        \midrule
        \textbf{2} & 1 & S  & $\emptyset$ & $\emptyset$ \\
                   & 2 & A  & [B, D]      & \{S\} \\
                   & 3 & B  & [D]         & \{A, S\} \\
                   & 4 & D  & $\emptyset$ & \{B, A, S\} \\
                   & 5 & C  & [E]         & \{B, D, A, S\} \\
                   & 6 & E  & $\emptyset$ & \{B, C, D, A, S\} \\
        \midrule
        \textbf{3} & 1 & S  & $\emptyset$ & $\emptyset$ \\
                   & 2 & A  & [B, D]      & \{S\} \\
                   & 3 & B  & [D]         & \{A, S\} \\
                   & 4 & C  & [D, G1]     & \{B, A, S\} \\
                   & 5 & G1 & [D]         & \{B, C, A, S\} \\
        \bottomrule
    \end{tabular}
    
    \vspace{0.2cm}
    \textbf{Cammino:} S $\rightarrow$ A $\rightarrow$ B $\rightarrow$ G1 \quad \textbf{Costo:} 15

\begin{center}
\begin{forest}
  for tree={circle, minimum size=1.5em, inner sep=2pt, s sep=4mm, l sep=5mm, unexplored}
  [S, explored
    [A, explored
      [B, explored
        [C, explored]
        [G1, explored]
      ]
    ]
    [B]
    [D, explored]
  ]
\end{forest}
\end{center}

\end{table}
\begin{itemize}
    \item  4 iterazioni di limite (da 0 a 3), 16 nodi espansi (1+4+6+5)
    \item il cammino trovato non è ottimo (costo 15 contro 14)
    \item come la BFS
\end{itemize}


% ---------------------------------------------------------
% Tabella Greedy 
% ---------------------------------------------------------
\begin{table}[H]
    \centering
    \caption{ricerca best-first greedy}
    \begin{tabular}{clp{4cm}p{5cm}}
        \toprule
        \textbf{Passo} & \textbf{Nodo Espanso} & \textbf{Frontiera} & \textbf{Esplorati} \\
        \midrule
        1 & S  & $\emptyset$ & $\emptyset$ \\
        2 & B  & [A, D]      & \{S\} \\
        3 & G1 & [A, D, C]   & \{B, S\} \\
        \bottomrule
    \end{tabular}
    
    \vspace{0.2cm}
    \textbf{Cammino:} S $\rightarrow$ B $\rightarrow$ G1 \quad \textbf{Costo:} 16

\begin{center}
\begin{forest}
  for tree={circle, minimum size=1.5em, inner sep=2pt, s sep=4mm, l sep=5mm, unexplored}
  [S, explored
    [A]
    [B, explored
      [C]
      [G1, explored]
    ]
    [D]
  ]
\end{forest}
\end{center}
\end{table}
\begin{itemize}
    \item  3 iterazioni
    \item  il cammino trovato non è ottimo $16 > 14$ 
    \item  greedy non garantisce l'ottimalità, poiché si affida esclusivamente all'euristica ignorando il costo del cammino percorso
\end{itemize}

% ---------------------------------------------------------
% Tabella A*
% ---------------------------------------------------------

\begin{table}[H]
    \centering
    \caption{ricerca A*}
    \begin{tabular}{clp{4cm}p{5cm}}
        \toprule
        \textbf{Passo} & \textbf{Nodo Espanso} & \textbf{Frontiera} & \textbf{Esplorati} \\
        \midrule
        1 & S  & $\emptyset$       & $\emptyset$ \\
        2 & D  & [A, B]            & \{S\} \\
        3 & B  & [A, C, E]         & \{D, S\} \\
        4 & C  & [A, E, G1]        & \{B, D, S\} \\
        5 & A  & [E, G1, F, J]     & \{B, C, D, S\} \\
        6 & E  & [G1, F, J]        & \{B, A, C, D, S\} \\
        7 & F  & [G1, J, G2]       & \{E, B, A, C, D, S\} \\
        8 & G2 & [G1, J]           & \{F, E, B, A, C, D, S\} \\
        \bottomrule
    \end{tabular}
    
    \vspace{0.2cm}
    \textbf{Cammino:} S $\rightarrow$ D $\rightarrow$ C $\rightarrow$ F $\rightarrow$ G2 \quad \textbf{Costo:} 14


\begin{center}
\begin{forest}
  for tree={circle, minimum size=1.5em, inner sep=2pt, s sep=4mm, l sep=5mm, unexplored}
  [S, explored
    [A, explored]
    [B, explored
      [G1]
    ]
    [D, explored
      [C, explored
        [F, explored
          [G2, explored]
        ]
        [J]
      ]
      [E, explored]
    ]
  ]
\end{forest}
\end{center}
\end{table}

\begin{itemize}
    \item sono state necessarie 8 iterazioni per trovare l'obiettivo
    \item il cammino trovato è ottimo (14)
    \item A* garantisce l'ottimalità se l'euristica è ammissibile (nella ricerca su alberi) o consistente (nella ricerca su grafi)
\end{itemize}

\subsection*{euristica}
L'euristica è:
\begin{itemize}
    \item \textbf{ammissibile} - perché per ogni stato $s$ vale $h(s) \le$ costo minimo per raggiungere l'obiettivo
    \item \textbf{non consistente} - perché $9 = h(A) > c(A, B) + h(B) = 4 + 3 = 7$
\end{itemize}


\end{document}
